\chapter{中心对称与平行四边形}\label{ch:g7:central}

将图形绕一点旋转半圈，即\emph{中心
对称}，是本书继反射之后的第二种变换（见
\cref{ch:g6:symmetry}）。其典型图形为平行四边形 ---
即关于自身两条对角线交点中心对称的
四边形。

\section{中心对称}

\begin{definition}[关于一点的对称]\label{def:g7:central:def}
点 $M$ 关于一点 $O$ 的\emph{对称点}\index{中心对称}
是点 $M'$，使 $O$ 为
$[MM']$ 的\emph{中点}。对称图形即原图形经
$180^\circ$ 的旋转（绕 $O$）后所得的图形。
\end{definition}

\begin{omfigure}
\begin{tikzpicture}[scale=1.0]
  \fill (0,0) circle (1.6pt) node[below, font=\small] {$O$};
  \fill[omDef] (-2.3,1.0) circle (1.8pt) node[above, font=\small] {$M$};
  \fill[omThm] (2.3,-1.0) circle (1.8pt) node[below, font=\small] {$M'$};
  \draw[dashed] (-2.3,1.0) -- (2.3,-1.0);
  \draw (-1.2,0.45) -- (-1.1,0.6);
  \draw (1.1,-0.6) -- (1.2,-0.45);
  \fill[omDef] (-1.4,-0.8) circle (1.8pt) node[below, font=\small] {$N$};
  \fill[omThm] (1.4,0.8) circle (1.8pt) node[above, font=\small] {$N'$};
  \draw[dashed] (-1.4,-0.8) -- (1.4,0.8);
\end{tikzpicture}

{\small 关于 $O$ 的对称：每一点及其像点都与
$O$ 共线，位于其两侧且等距（图中刻线标出）。}
\end{omfigure}

\begin{method}[作一点的对称点]\label{met:g7:central:construct}
求作 $M'$，它是 $M$ 关于 $O$ 的对称点，作法如下：
\begin{enumerate}
  \item 作直线 $(MO)$，并将其向 $O$ 之外延长；
  \item 用直尺或圆规量出 $MO$；
  \item 在延长线上取点 $M'$，使 $OM' = OM$；
\end{enumerate}
在方格纸上：数出从 $M$ 到 $O$ 的水平与竖直步数，
再从 $O$ 出发走\emph{相同}的步数到达 $M'$。
\end{method}

\begin{example}[在方格纸上]\label{ex:g7:central:grid}
若 $M$ 在左 $3$ 格、上 $1$ 格处（以 $O$ 为基准），其对称点
$M'$ 位于 \emph{右} $3$ 格、\emph{下} $1$ 格处（以 $O$ 为基准）：
中心对称同时改变两个方向（反射只改变
一个方向）。
\end{example}

\begin{proposition}[半圈旋转保持的性质]\label{prop:g7:central:preserved}
中心对称保持长度、角、周长与面积；
直线的像是与之\emph{平行}的直线，线段的像是
与之平行且等长的线段，圆的像是
同半径的圆（圆心为原圆心的对称点）。
\end{proposition}
\admitted

\begin{remark}
与反射对比（见\cref{prop:g6:symmetry:preserved}）：两者都保持
长度与角；但反射会将图形翻转过来
（如“b”变成“d”），而半圈旋转使图形保持可读 ---
即上下颠倒（如“b”变成“q”）。与反射不同，
半圈旋转使每条直线都变为与之平行的直线。
\end{remark}

\begin{definition}[对称中心]\label{def:g7:central:center}
若将图形绕 $O$ 旋转半圈后恰好与自身重合，
则点 $O$ 称为该图形的\emph{对称中心}\index{对称中心}。
例如字母 S、N、Z 都有对称中心；
圆有一个对称中心（即圆心）；三角形则永远没有对称中心。
\end{definition}

\section{平行四边形}

\begin{definition}[平行四边形]\label{def:g7:central:parallelogram}
若四边形的两条对角线相交于公共中点，
则称该四边形为\emph{平行四边形}\index{平行四边形}。
此交点即为图形的对称中心：每个顶点都是
其对顶点经半圈旋转所得的像。
\end{definition}

\begin{theorem}[平行四边形的性质]\label{thm:g7:central:properties}
在平行四边形 $ABCD$ 中，中心为 $O$：
\begin{enumerate}
  \item 对边平行：$(AB) \parallel (CD)$ 且
        $(BC) \parallel (AD)$；
  \item 对边相等：$AB = CD$ 且
        $BC = AD$；
  \item 对角相等：$\widehat A = \widehat C$ 且
        $\widehat B = \widehat D$。
\end{enumerate}
\end{theorem}

\begin{proof}
绕点 $O$ 的半圈旋转把 $A$ 送到 $C$，$B$ 送到 $D$（由定义：
$O$ 为两条对角线的中点）。于是它把线段
$[AB]$ 映为线段 $[CD]$；由
\cref{prop:g7:central:preserved}可知，它的像与 $[AB]$ 平行且
等长，即结论 1、2 成立。它也把 $A$ 处的角
送到 $C$ 处的角，角保持不变，即结论 3 成立。
\end{proof}

\begin{omfigure}
\begin{tikzpicture}[scale=1.0]
  \coordinate (A) at (0,0);
  \coordinate (B) at (3.4,0);
  \coordinate (C) at (4.6,1.9);
  \coordinate (D) at (1.2,1.9);
  \draw[thick] (A) -- (B) -- (C) -- (D) -- cycle;
  \draw[omThm] (A) -- (C);
  \draw[omThm] (B) -- (D);
  \fill (2.3,0.95) circle (1.6pt) node[right=3pt, font=\small] {$O$};
  \node[below left, font=\small] at (A) {$A$};
  \node[below right, font=\small] at (B) {$B$};
  \node[above right, font=\small] at (C) {$C$};
  \node[above left, font=\small] at (D) {$D$};
  \draw (1.62,-0.09) -- (1.78,0.09);
  \draw (2.82,1.81) -- (2.98,1.99);
  \draw (0.53,1.02) -- (0.67,0.88);
  \draw (3.93,1.02) -- (4.07,0.88);
\end{tikzpicture}

{\small 平行四边形及其对称中心 $O$：对角线（红色）
彼此相交于各自的中点，对边平行且相等
（刻线标记）。}
\end{omfigure}

\begin{theorem}[判定平行四边形]\label{thm:g7:central:recognize}
只要四边形（各边不相交）满足
以下条件\emph{之一}，即为平行四边形：
\begin{enumerate}
  \item 对角线有相同的中点（即定义）；
  \item 两组对边分别平行；
  \item 一组对边平行，\emph{且}长度相等。
\end{enumerate}
\end{theorem}
\admitted

\begin{example}\label{ex:g7:central:recognize}
$ABCD$ 满足 $(AB) \parallel (CD)$，且 $AB = CD = 5$ cm。由判定条件 3
可知 $ABCD$ 是平行四边形 --- 因此无需再测量，
便可得 $AD = BC$，$(AD) \parallel (BC)$，以及
其对角线互相平分。
\end{example}

\begin{remark}[特殊的平行四边形]
\cref{ch:g6:shapes}中的长方形、菱形与正方形，恰好是
具有附加性质的平行四边形：对角线相等
（长方形）、对角线垂直（菱形），两者兼具（正方形）。
\cref{ex:g6:shapes:recognize}中许诺的角度论证现在行得通：
在平行四边形中，邻角互补
（借助平行线的内错角与同位角，
\cref{thm:g7:triangles:equalities}），故一个直角即可推出
四个角都是直角。
\end{remark}

\section{练习}

\begin{exercise}[$\star$]\label{exo:g7:central:1}
在方格纸上标出 $O$，再标出点 $M$，它在 $O$ 右侧
四格、上方一格处。作出 $M'$，即 $M$ 关于 $O$ 的对称点，以及
$M$ 关于过 $O$ 的\emph{竖直方格线}的对称点。
比较这两个像。
\end{exercise}

\begin{exercise}[$\star$]\label{exo:g7:central:2}
画三角形 $ABC$ 与三角形外一点 $O$，作出它的像
$A'B'C'$，即绕 $O$ 作半转所得的图形。
关于 $A'B'$ 与 $AB$ 的长度，以及直线 $(A'B')$ 与 $(AB)$，你能说明什么？
\end{exercise}

\begin{exercise}[$\star$]\label{exo:g7:central:3}
哪些数字（按计算器显示屏上的写法，$0$ 到 $9$）具有
对称中心？大写字母 H、A、N、S、T、Z、O 中哪些具有对称中心？
\end{exercise}

\begin{exercise}[$\star$]\label{exo:g7:central:4}
$KLMN$ 是对称中心为 $O$ 的平行四边形，$KL = 7$ cm，
$LM = 4$ cm，$\widehat K = 115^\circ$。求下列量，并说明理由：$MN$、
$NK$、$\widehat M$ 以及 $[LN]$ 的中点。
\end{exercise}

\begin{exercise}[$\star$]\label{exo:g7:central:5}
由对角线作平行四边形 $ABCD$：$AC = 8$ cm，
$BD = 5$ cm，两条对角线在各自的中点处相交，夹角为 $60^\circ$。
（先画出对角线！）
\end{exercise}

\begin{exercise}[$\star$]\label{exo:g7:central:6}
在平行四边形中，一个角为 $115^\circ$。
利用关于邻角互补的注，求其余三个角。
\end{exercise}

\begin{exercise}[$\star\star$]\label{exo:g7:central:7}
描出 $A(1, 1)$、$B(4, 2)$ 及点 $O(2.5, 2.5)$。计算
像 $A'$ 与 $B'$ 的坐标（绕 $O$ 作半转）
（利用\cref{ex:g7:central:grid}中的数格方法）。
\end{exercise}

\begin{exercise}[$\star\star$]\label{exo:g7:central:8}
$ABCD$ 是四边形，其中 $AB = CD$。仅凭此条件能否使其
成为平行四边形？若不能，画出一个反例（等腰梯形有助于思考），
并指出应在假设中补充什么条件。
\end{exercise}

\begin{exercise}[$\star\star$]\label{exo:g7:central:9}
设 $ABC$ 为三角形，$O$ 为 $[BC]$ 的中点。作出
点 $A'$，即 $A$ 关于 $O$ 的对称点。证明 $ABA'C$ 是
平行四边形。（\cref{thm:g7:central:recognize}中的哪一条判定在这里
无需额外验证即可满足？）
\end{exercise}

\begin{exercise}[$\star\star$]\label{exo:g7:central:10}
判断下列说法是否正确，并说明理由：“对角两两相等的四边形
是平行四边形”；“对角线互相垂直的平行四边形
是菱形”；“有一个直角的平行四边形
是长方形”。
\end{exercise}

\begin{exercise}[$\star\star\star$]\label{exo:g7:central:11}
设 $ABCD$ 为任意四边形，$I$、$J$、$K$、$L$ 分别是
其边 $[AB]$、$[BC]$、$[CD]$、$[DA]$ 的中点。画出若干例子。
$IJKL$ 总是什么图形？（证明借助中点
定理，将在\cref{ch:g8:midpoints}中给出 --- 并在
高中分册中再借助向量证明。）
\end{exercise}

\section{问题：蛋糕、中线与不可能存在的第二个
对称中心}

\begin{problem}[{周末问题 --- 半转的实际应用：公平切蛋糕、
中线的界，以及为何没有图形能有两个
对称中心}]
\label{pb:g7:central:1}
半转看似最简单的变换 --- 然而
它能把蛋糕切成可证明公平的两半，度量
三角形的中线，并蕴含一颗纯数学的小珍珠：一个
有界图形可以有一个对称中心，但\emph{绝不
可能有两个}。证明用到你将在第 I 部分中发现的一个结论：两个
半转依次作用，合成为一个平移。

\medskip
\noindent\textbf{第 I 部分 --- 半转体操。}
\begin{enumerate}
  \item 在方格纸上，取原点 $O(0, 0)$ 为中心。
        利用\cref{ex:g7:central:grid}中的方格计数，
        求 $M(3, 1)$、$N(-2, 4)$ 和
        $P(0, -5)$ 关于 $O$ 的半转下的像。$(x, y)$ 的
        一般规律是什么？
  \item 现取中心为 $C(2, 1)$，求
        $M(5, 3)$ 的像，并检验一般公式：像的每个坐标
        都是\emph{中心坐标的两倍减去该点坐标}
        的结果。
  \item 扑克牌的图案设计成倒过来看也一样，
        ——这是一种中心对称，因此两位玩家都不会
        “看错”牌面。试解释：在点数为\emph{奇数}的牌上，
        必有一个点恰好位于牌的
        中心。
  \item 连续两次半转：取中心 $O_1(2, 0)$
        与 $O_2(5, 0)$，将点 $A(1, 1)$ 先经关于 $O_1$ 的半转
        变换，再将所得的像经关于 $O_2$ 的半转。
        比较起点与终点。对 $B(0, 3)$ 重复上述过程。
        两次半转合起来相当于哪一种简单的
        单一变换？
  \item 将第 4 题中的平移量与距离
        $O_1 O_2$ 作比较，并叙述这一发现。（与
        \cref{pb:g6:symmetry:1}中两面平行镜子的情形作比较：
        现象相同，只是主角换成了新的。）
\end{enumerate}

\medskip
\noindent\textbf{第 II 部分 --- 蛋糕与中线。}
\begin{enumerate}[resume]
  \item 设 $O$ 为平行四边形的中心。试解释为何
        过 $O$ 的任一直线与边界交于两点，
        这两点关于 $O$ 对称，并解释为何该直线
        将平行四边形分成两块，它们恰好是
        彼此经半转所得的全等图形，因而面积相等。
  \item 面包师定理：长方形蛋糕可被
        \emph{任一}过其中心的直线切成完全公平的两份。更进一步，
        长方形蛋糕上可以挖出一个长方形的洞（大小、位置任意，
        甚至可以倾斜）。试描述一条直线切法，
        使切出的两部分蛋糕分量相等，
        并利用第 6 题加以说明。
  \item 平行四边形 $ABCD$ 的三个顶点为 $A(0, 0)$、
        $B(6, 1)$ 和 $C(8, 5)$。求顶点 $D$ 的坐标
        以及中心 $O$ 的坐标。
  \item \Cref{exo:g7:central:9}中，由三角形 $ABC$ 出发，
        设 $O$ 为 $[BC]$ 的中点，作出了平行四边形
        $ABA'C$，其中 $A'$ 为 $A$ 关于 $O$ 的对称点。
        利用它，结合三角形不等式
        （\cref{thm:g7:triangles:inequality}），在三角形
        $ABA'$ 中证明\emph{中线界}：中线
        $[AO]$ 满足下式：
        \[
        AO < \frac{AB + AC}{2} .
        \]
        （注意：$AA' = 2\,AO$，$BA' = AC$。）
  \item 在三角形中，$AB = 5$ 厘米，$AC = 7$ 厘米，
        由 $A$ 引出的中线长度应介于哪两个数值之间？
        （请同时利用三角形 $ABA'$ 中
        不等式的两个方向。）
\end{enumerate}

\medskip
\noindent\textbf{第 III 部分 --- 一个图形可以有几个
对称中心？}
\begin{enumerate}[resume]
  \item 对下列每个图形，判断它是否有对称中心，
        若有，指出其位置：线段；整条直线；
        等边三角形；正方形；圆；
        字母 S 与 Z。
  \item 证明你对等边三角形的结论：若一个半转将三角形
        映回自身，则它会把三个顶点两两配对
        ——但三是奇数，因此必有某个顶点
        被映到自身。这将迫使该顶点成为什么？
        为何这是不可能的？
  \item 推广第 12 题：解释为何\emph{没有}任何顶点数为
        奇数的多边形具有
        对称中心。
  \item 珍珠问题：假设某图形有\emph{两个}不同的
        对称中心 $O_1$ 与 $O_2$。依次施行这两次半转，
        并利用第 5 题的发现：
        哪一种运动必须将该图形恰好映回自身？
        试解释为何能放在一张纸上的图形
        无法满足这一点，并由此得出结论。
  \item 第 14 题的结论留下了一个漏洞：
        \emph{无界}图形。在一条无限饰带上验证，
        ——例如一串无尽的脚印，依次为左、右、
        左、右\,\dots\ ——它（至少）有两个
        不同的对称中心，并指出第 14 题所预言的
        那个平移。
\end{enumerate}
\end{problem}
